\documentclass[11pt]{amsart}

\usepackage[T1]{fontenc}
\usepackage{lmodern}
\usepackage{microtype}
\usepackage{amsmath,amssymb,amsthm}
\usepackage[hidelinks]{hyperref}

\newtheorem{theorem}{Theorem}
\theoremstyle{remark}
\newtheorem{remark}[theorem]{Remark}

\newcommand{\disc}{\operatorname{disc}}
\newcommand{\dc}{\operatorname{dc}}

\title[A rank-zero counterexample]
{A Rank-Zero Counterexample to Conjecture~8.4 of
Abdi--Cornu\'ejols--Zlatin}

\date{}

\subjclass[2020]{05C20, 05C70}
\keywords{dicut, dijoin, rounded $1$-factor, cubic bipartite multigraph,
Hamiltonian cycle}

\hypersetup{
  pdftitle={
    A Rank-Zero Counterexample to Conjecture 8.4 of
    Abdi-Cornuejols-Zlatin
  }
}

\begin{document}

\begin{abstract}
We disprove Conjecture~8.4 of Abdi, Cornu\'ejols, and Zlatin as
stated. Orienting the eight-vertex planar cubic bipartite multigraph
$1110$ of Exoo and Harary from one bipartition class to the other
produces a planar sink-regular $(3,4)$-bipartite digraph $D$ with
$\rho(3,D)=0$. The matroid $M_1(D,\mathbf{1})$ has unique basis
$\varnothing$, while every rounded $1$-factor has disconnected
complement.
\end{abstract}

\maketitle

Conjecture~8.4 of Abdi, Cornu\'ejols, and Zlatin asserts that, for a
planar sink-regular $(3,4)$-bipartite digraph $D=(V,A)$ and disjoint
bases $Q_1,Q_2,Q_3$ of $M_1(D,\mathbf{1})$, there is a rounded
$1$-factor $J$ such that $\dc(J)=Q_1$ and $D\setminus J$ is connected
\cite[Conjecture~8.4]{ACZ}. Their conventions permit parallel arcs,
and planarity and connectedness refer to the underlying undirected
multigraph \cite[pp.~2418, 2427]{ACZ}.

\begin{theorem}
Conjecture~8.4 of \cite{ACZ} is false, already for an eight-vertex
planar instance with $\rho(3,D)=0$.
\end{theorem}

\begin{proof}
Let
\[
 S=\{u,b_1,b_2,b_3\},
 \qquad
 T=\{v,a_1,a_2,a_3\}.
\]
For each $i\in\{1,2,3\}$, add one edge $ua_i$, two parallel edges
between $a_i$ and $b_i$, and one edge $b_iv$. This is the multigraph
$1110$ in \cite[Figure~3]{EH}. Drawing its three $u$--$v$ strands in
parallel shows that the resulting graph $G$ is planar; it is also
connected, cubic, and bipartite.

Orient every edge from $S$ to $T$, obtaining $D$. Let
$\delta_D^+(U)$ be a dicut, and put $X=U\cap S$ and $Y=U\cap T$.
Since $\delta_D^-(U)=\varnothing$,
\[
 |\delta_D^+(U)|
 =\sum_{z\in U}\bigl(d_D^+(z)-d_D^-(z)\bigr)
 =3|X|-3|Y|.
\]
The underlying graph is connected and $U$ is nonempty and proper, so
the cut is nonempty. Thus every dicut has size at least $3$. Since
all vertices have degree $3$, the digraph $D$ is planar and
sink-regular $(3,4)$-bipartite.

There are no active vertices, and
\[
 \disc(V)=|T|-|S|=0,
 \qquad
 \rho(3,D)=0
\]
by \cite[Lemma~3.7(1)]{ACZ}. Moreover, for every dicut shore $U$,
\[
 \disc(U)=|Y|-|X|
 =-\frac{1}{3}|\delta_D^+(U)|\le -1.
\]
By \cite[Definition~4.7]{ACZ}, $\varnothing$ is therefore the unique
basis of $M_1(D,\mathbf{1})$: it has cardinality $\disc(V)=0$ and
satisfies
\[
 0=|\varnothing\cap U|\ge 1+\disc(U)
\]
for every dicut shore. Set
\[
 Q_1=Q_2=Q_3=\varnothing.
\]
Each $Q_i$ is a basis, and $Q_i\cap Q_j=\varnothing$ whenever
$i\ne j$.

Let $J$ be a rounded $1$-factor. Since every vertex has degree $3$,
Definition~3.3 of \cite{ACZ} gives $|J\cap\delta_D(z)|=1$ at every
vertex $z$. Hence $J$ is a perfect matching and
$\dc(J)=\varnothing=Q_1$.

There is a unique $i$ for which $ua_i\in J$. The vertex $a_i$ is
then matched to $u$, so neither parallel edge between $a_i$ and $b_i$
belongs to $J$. Consequently $b_iv\in J$. After deleting $J$, the
two parallel edges between $a_i$ and $b_i$ form a connected component
of the underlying graph of $D\setminus J$. Hence $D\setminus J$ is
disconnected. No rounded $1$-factor satisfies the conjectured
conclusion.
\end{proof}

\begin{remark}[Scope]
The example lies in the rank-zero case and is not $3$-connected:
removing $u$ and $v$ leaves the three components $\{a_i,b_i\}$.
Thus it does not address a positive-rank reformulation of
Conjecture~8.4 or Barnette's conjecture. Nor does it disprove the
planar three-dijoin statement in \cite[Theorem~8.5]{ACZ}; it only
falsifies Conjecture~8.4 as published.

More generally, the dicut and matroid calculations above apply to any
connected cubic bipartite graph oriented from one bipartition class to
the other, and a connected complement of a rounded $1$-factor would
be a Hamiltonian circuit. Thus simplicity alone does not repair the
statement: the $26$-vertex simple planar cubic bipartite
non-Hamiltonian graph of Asano, Saito, Exoo, and Harary gives a simple
counterexample \cite{ASEH}.
\end{remark}

\begin{thebibliography}{9}

\bibitem{ACZ}
A.~Abdi, G.~Cornu\'ejols, and M.~Zlatin,
\newblock \emph{On packing dijoins in digraphs and weighted digraphs},
\newblock SIAM J. Discrete Math. \textbf{37} (2023), no.~4,
2417--2461,
\href{https://doi.org/10.1137/22M1506511}
{doi:10.1137/22M1506511}.

\bibitem{ASEH}
T.~Asano, N.~Saito, G.~Exoo, and F.~Harary,
\newblock \emph{The smallest $2$-connected cubic bipartite planar
nonhamiltonian graph},
\newblock Discrete Math. \textbf{38} (1982), no.~1, 1--6,
\href{https://doi.org/10.1016/0012-365X(82)90162-5}
{doi:10.1016/0012-365X(82)90162-5}.

\bibitem{EH}
G.~Exoo and F.~Harary,
\newblock \emph{The smallest cubic multigraphs with prescribed
bipartite, block, Hamiltonian and planar properties},
\newblock Indian J. Pure Appl. Math. \textbf{13} (1982), no.~3,
309--312.

\end{thebibliography}

\end{document}